\section{Adaptive interaction graphs}
\subsection{A local graph with a controlled expansion}
For particle $i$, let $x_i^t\in\mathbb R^d$ denote position. The simulator receives six consecutive frames $H_t$, particle attributes and boundary information, and predicts normalized acceleration $\mu_i^t$ and residual variance $q_i^t>0$. It reverses the training normalization before integrating acceleration to update positions.

Let $E_0(x_t)$ be the simulator's native directed base graph, including self-messages and its 128-neighbor receiver cap. We preserve that graph and construct optional unordered pairs in a surrounding annulus:
\begin{equation}
 C_R(x_t)=\{\{i,j\}:i<j,\ r<\|x_i^t-x_j^t\|\le R\}.
\end{equation}
Short-range pairs omitted by the native cap do not enter the annulus. A policy selects $S_t\subseteq C_R(x_t)$ and appends both orientations of each selected pair. For an integer budget $B_t\ge0$,
\begin{equation}
 |E_t|=|E_0(x_t)|+2\min(B_t,|C_R(x_t)|).
 \label{eq:budget}
\end{equation}
Our budgeted policies use $B_t=\lfloor0.25|C_R(x_t)|\rfloor$. They therefore add identical numbers of edges on a common particle state, while differing in placement. The budget controls additions, not total cost: the base graph grows as particles cluster, and appended pairs are not recapped, so final receiver degrees may exceed 128. Figure~\ref{fig:interaction-schematic} illustrates the construction.

\subsection{Learning to use additional messages}
A model trained only on its native graph need not benefit from extra messages. We expose a paired training arm to the same kind of graph changes used at evaluation. Independently for each training example, with probability $1/2$ we append a uniformly random quarter of the annulus pairs. The base-only arm retains its native graph. Both arms share the faithful objective, initial weights, sampled frames, noise and optimization schedule. Holding the evaluation policy fixed then measures the effect of this graph-exposure intervention. The residual head still learns to predict error; exposure does not give it an action-benefit target.

\subsection{Allocating with cached residual scores}
Residual allocation scores each optional pair by its larger endpoint score,
\begin{equation}
 s_{t,\{i,j\}}=\max(q_i^{t-1},q_j^{t-1}),
\end{equation}
and selects the highest-scoring $\min(B_t,|C_R|)$ pairs. A single base-graph pass on the initial six observed frames initializes the cache. Thereafter, each forecast follows the same cycle: select current pairs using cached scores, run the simulator, update positions, and cache its new scores. The budget applies from forecast 1. No future positions or labels enter the rule, and no separate scoring pass is needed at subsequent forecasts. Appendix~\ref{sec:implementation-details} specifies tie handling and the compact pilot's initialization.

This autonomous controller must be distinguished from the observed-history diagnostic. There, scores come from the preceding observed history on the base graph, and all policies receive the same current history. During a rollout, scores instead come from the controller's own preceding selected graph. The first comparison isolates placement; the second includes feedback.

\begin{figure}[t]
\centering\setlength{\unitlength}{1pt}
\begin{picture}(225,105)
\put(35,96){\makebox(0,0){\scriptsize Local graph}}
\linethickness{0.35pt}
\qbezier(11,38)(14.5,46)(18,54)
\linethickness{0.35pt}
\qbezier(34,33)(36,41.5)(38,50)
\linethickness{0.35pt}
\qbezier(38,50)(45.5,55.5)(53,61)
\linethickness{0.35pt}
\qbezier(53,61)(56.5,52)(60,43)
\put(11,38){\circle*{3}}
\put(18,54){\circle*{3}}
\put(34,33){\circle*{3}}
\put(38,50){\circle*{3}}
\put(53,61){\circle*{3}}
\put(60,43){\circle*{3}}
\put(22,76){\circle*{3}}
\put(44,79){\circle*{3}}
\put(35,13){\makebox(0,0){\tiny 4 mandatory}}
\put(111,96){\makebox(0,0){\scriptsize Candidate pairs}}
\linethickness{0.35pt}
\qbezier(87,38)(90.5,46)(94,54)
\linethickness{0.35pt}
\qbezier(110,33)(112,41.5)(114,50)
\linethickness{0.35pt}
\qbezier(114,50)(121.5,55.5)(129,61)
\linethickness{0.35pt}
\qbezier(129,61)(132.5,52)(136,43)
\linethickness{0.25pt}
\qbezier[10](87,38)(98.5,35.5)(110,33)
\linethickness{0.25pt}
\qbezier[10](94,54)(104,52)(114,50)
\linethickness{0.25pt}
\qbezier[10](94,54)(96,65)(98,76)
\linethickness{0.25pt}
\qbezier[10](114,50)(125,46.5)(136,43)
\linethickness{0.25pt}
\qbezier[10](129,61)(124.5,70)(120,79)
\linethickness{0.25pt}
\qbezier[10](98,76)(109,77.5)(120,79)
\put(87,38){\circle*{3}}
\put(94,54){\circle*{3}}
\put(110,33){\circle*{3}}
\put(114,50){\circle*{3}}
\put(129,61){\circle*{3}}
\put(136,43){\circle*{3}}
\put(98,76){\circle*{3}}
\put(120,79){\circle*{3}}
\put(111,13){\makebox(0,0){\tiny 6 optional}}
\put(187,96){\makebox(0,0){\scriptsize Budgeted graph}}
\linethickness{0.35pt}
\qbezier(163,38)(166.5,46)(170,54)
\linethickness{0.35pt}
\qbezier(186,33)(188,41.5)(190,50)
\linethickness{0.35pt}
\qbezier(190,50)(197.5,55.5)(205,61)
\linethickness{0.35pt}
\qbezier(205,61)(208.5,52)(212,43)
\linethickness{0.9pt}
\qbezier(163,38)(174.5,35.5)(186,33)
\put(163,38){\circle*{3}}
\put(170,54){\circle*{3}}
\put(186,33){\circle*{3}}
\put(186,33){\circle{6}}
\put(190,50){\circle*{3}}
\put(190,50){\circle{6}}
\put(205,61){\circle*{3}}
\put(212,43){\circle*{3}}
\put(174,76){\circle*{3}}
\put(196,79){\circle*{3}}
\put(187,13){\makebox(0,0){\tiny 4 mandatory + 1 added}}
\end{picture}
\caption{A controlled graph update. The mandatory local graph is retained, and the budget adds one of six candidate annulus pairs. Residual allocation prioritizes cached endpoint scores (circled particles have the largest scores). The thick edge is the selected addition; each optional pair enters the simulator in both directions.}
\label{fig:interaction-schematic}
\end{figure}

\subsection{Learning the residual signal without changing mean fitting}
We use established faithful heteroscedastic regression \citep{stirn2023} to keep residual learning from reweighting the mean predictor's error. Let $h_i$ be the shared representation, $y_i$ the normalized target, $R_i=\|\mu_i-y_i\|_2^2$, and $\operatorname{sg}$ denote stop-gradient. Ordinary isotropic Gaussian negative log-likelihood, omitting its constant, is
\begin{align}
 \ell_{\mathrm{NLL},i}&=\frac{R_i}{2q_i}+\frac d2\log q_i,\nonumber\\
 q_i&=\operatorname{softplus}(g_\phi(h_i))+q_{\min}.
 \label{eq:nll}
\end{align}
The faithful variant detaches both the shared features seen by the variance head and its residual target:
\begin{align}
 q_i&=\operatorname{softplus}(g_\phi(\operatorname{sg}(h_i)))+q_{\min},\nonumber\\
 \ell_{\mathrm F,i}&=\tfrac12R_i+
 \frac{\operatorname{sg}(R_i)}{2q_i}+\frac d2\log q_i.
 \label{eq:faithful}
\end{align}
Only the first term updates the mean network, yielding the same raw mean gradient as squared-error training on matched inputs and graphs. Separate gradient clipping preserves this separation. The scalar $q_i$ predicts per-coordinate residual variance in normalized acceleration units; $dq_i$ predicts vector squared error. It is not an isolated measure of epistemic uncertainty or physical stochasticity. Objective and normalization details appear in Appendix~\ref{sec:implementation-details}.
